\documentclass[11pt,a4paper]{article}

\usepackage[margin=2.5cm]{geometry}
\usepackage[T1]{fontenc}
\usepackage{lmodern}
\usepackage{microtype}
\usepackage{amsmath,amssymb,bm}
\usepackage{booktabs}
\usepackage{xcolor}
\usepackage{hyperref}
\usepackage{listings}

\hypersetup{
  colorlinks=true,
  linkcolor=blue!55!black,
  urlcolor=blue!55!black,
  citecolor=blue!55!black
}

\lstdefinestyle{pythonstyle}{
  language=Python,
  basicstyle=\ttfamily\small,
  keywordstyle=\color{blue!65!black},
  commentstyle=\color{green!40!black},
  stringstyle=\color{orange!60!black},
  numbers=left,
  numberstyle=\tiny\color{gray},
  frame=single,
  breaklines=true,
  showstringspaces=false,
  columns=fullflexible
}

\newcommand{\vect}[1]{\bm{#1}}
\newcommand{\ones}{\vect{1}}

\title{Implementation Note: Direct Gaussian QEq and External-Field Regularization}
\author{SchnetPack charge-prediction implementation}
\date{\today}

\begin{document}
\maketitle

\begin{abstract}
This note relates the implementation of \texttt{EMLEQEqStatic} in
\path{src/schnetpack/atomistic/atomwise.py} to the underlying constrained
charge-equilibration (QEq) equations.  The implementation constructs a dense
Gaussian Coulomb matrix and solves the charge constraint directly with
\texttt{torch.linalg.solve}.  It computes both a vacuum charge distribution and,
when an external electrostatic potential is supplied, a polarized charge
distribution.  The appendix documents the auxiliary field-regularized training
task implemented in \path{src/schnetpack/task.py} and configured in
\path{train_field_regularized.yaml}.
\end{abstract}

\section{Notation and units}

For one structure containing $N$ atoms, let
\begin{itemize}
  \item $q_i$ be the total charge on atom $i$, in units of the elementary charge $e$;
  \item $\chi_i$ be its environment-dependent electronegativity-like coefficient,
        in \(\mathrm{kJ\,mol^{-1}\,e^{-1}}\);
  \item $\phi_i$ be the external electrostatic potential at atom $i$, in
        \(\mathrm{kJ\,mol^{-1}\,e^{-1}}\);
  \item $\sigma_i$ be the Gaussian charge width, in \AA;
  \item $Q$ be the prescribed total charge of the structure, in $e$; and
  \item $A$ be the QEq interaction matrix, in
        \(\mathrm{kJ\,mol^{-1}\,e^{-2}}\).
\end{itemize}
Vectors are denoted by bold symbols, and $\ones$ is the vector containing $N$
ones.  The implementation uses a Coulomb conversion constant
$k_e=\texttt{K\_E\_KJMOL\_ANG\_E2}$ so that distances in \AA{} produce energies in
the units above.

\section{The quadratic QEq problem}

In vacuum, the charge-dependent energy is
\begin{equation}
  E_0(\vect q)
  = \vect\chi^{\mathsf T}\vect q
  + \frac{1}{2}\vect q^{\mathsf T}A\vect q,
  \label{eq:vacuum-energy}
\end{equation}
subject to the total-charge constraint
\begin{equation}
  \ones^{\mathsf T}\vect q=Q.
  \label{eq:charge-constraint}
\end{equation}
The linear term favors transfer according to the learned electronegativities,
whereas the quadratic term penalizes self-charging and describes interatomic
Coulomb interactions.

If an external potential is present, its coupling to the charges is
\begin{equation}
  E_{\mathrm{ext}}(\vect q)=\vect\phi^{\mathsf T}\vect q.
\end{equation}
The polarized QEq energy therefore becomes
\begin{equation}
  E_{\phi}(\vect q)
  =(\vect\chi+\vect\phi)^{\mathsf T}\vect q
  +\frac{1}{2}\vect q^{\mathsf T}A\vect q,
  \label{eq:polarized-energy}
\end{equation}
with the same constraint in Eq.~\eqref{eq:charge-constraint}.  The sign follows
the implementation convention: the external potential is added to $\vect\chi$
before solving.

\section{Learned parameters and Gaussian widths}

The neural-network parameter head produces an atom-wise valence width $s_i$ and
an electronegativity $\chi_i$.  A learnable global scale $a_{\mathrm{QEq}}$
converts the former into the QEq width,
\begin{equation}
  \sigma_i=a_{\mathrm{QEq}}s_i.
\end{equation}
The implementation is, schematically,
\begin{lstlisting}[style=pythonstyle]
s = inputs[self.widths_key]
chi = inputs[self.chi_key]
a_qeq = self._get_a_qeq()
sigma = a_qeq * s
\end{lstlisting}
Optional lower and upper bounds are then applied to $\sigma_i$.  These bounds
are numerical and physical robustness controls: a very large width yields a
small self-hardness and can make the model excessively polarizable, whereas a
vanishing width produces a very large hardness.

\section{Construction of the QEq matrix}

\subsection{Diagonal self-hardness}

For the Gaussian convention used here, the diagonal term is
\begin{equation}
  A_{ii}=J_{ii}=\frac{k_e}{\sqrt{\pi}\,\sigma_i}.
  \label{eq:self-hardness}
\end{equation}
The corresponding code is
\begin{lstlisting}[style=pythonstyle]
Jii = K_E_KJMOL_ANG_E2 / (sigma * math.sqrt(math.pi))
if jii_floor > 0.0:
    Jii = Jii + jii_floor
\end{lstlisting}
An optional positive \texttt{jii\_floor} increases every diagonal entry and
damps the response to an external field.

\subsection{Off-diagonal Gaussian Coulomb interaction}

For atoms $i\ne j$, define
\begin{align}
  R_{ij} &= \lVert\vect R_i-\vect R_j\rVert,\\
  \sigma_{ij} &= \sqrt{\sigma_i^2+\sigma_j^2}.
\end{align}
The off-diagonal matrix element is
\begin{equation}
  A_{ij}=J_{ij}
  =k_e\frac{\operatorname{erf}\!\left(
      R_{ij}/(\sqrt{2}\sigma_{ij})\right)}{R_{ij}}.
  \label{eq:gaussian-coulomb}
\end{equation}
This is implemented by
\begin{lstlisting}[style=pythonstyle]
dR = R[:, None, :] - R[None, :, :]
Rij = torch.linalg.norm(dR, dim=-1)
sigma_ij = torch.sqrt(
    sigma[:, None] ** 2 + sigma[None, :] ** 2
).clamp_min(self.eps_sigma)

Jij = (
    K_E_KJMOL_ANG_E2
    * torch.erf(Rij_safe / (math.sqrt(2.0) * sigma_ij))
    / Rij_safe
)
Jij = Jij.masked_fill(eye, 0.0)
A = torch.diag(Jii) + Jij
\end{lstlisting}
Gaussian smearing keeps the interaction finite at short range.  At large
separation, $\operatorname{erf}(x)\rightarrow1$, so Eq.~\eqref{eq:gaussian-coulomb}
approaches the point-charge interaction $k_e/R_{ij}$.  If
\texttt{offdiag\_cutoff} is set, off-diagonal terms beyond that distance are set
to zero.  The present training configuration uses a cutoff of $10$~\AA.

\section{Constrained direct solution}

Introduce a Lagrange multiplier $\lambda$ and the Lagrangian
\begin{equation}
  \mathcal L(\vect q,\lambda)
  =\vect g^{\mathsf T}\vect q
  +\frac{1}{2}\vect q^{\mathsf T}A\vect q
  +\lambda(\ones^{\mathsf T}\vect q-Q),
\end{equation}
where $\vect g=\vect\chi$ in vacuum and
$\vect g=\vect\chi+\vect\phi$ in the external field.  Stationarity gives
\begin{align}
  A\vect q+\vect g+\lambda\ones&=0,\\
  \ones^{\mathsf T}\vect q&=Q.
\end{align}
These equations form the augmented system
\begin{equation}
  \underbrace{
  \begin{bmatrix}
    A & \ones\\
    \ones^{\mathsf T} & 0
  \end{bmatrix}}_{K}
  \begin{bmatrix}
    \vect q\\ \lambda
  \end{bmatrix}
  =
  \begin{bmatrix}
    -\vect g\\ Q
  \end{bmatrix}.
  \label{eq:augmented-system}
\end{equation}
The Python implementation constructs Eq.~\eqref{eq:augmented-system} literally:
\begin{lstlisting}[style=pythonstyle]
n = A.shape[0]
K = torch.zeros((n + 1, n + 1), device=A.device, dtype=A.dtype)
K[:n, :n] = A
K[:n, n] = 1.0
K[n, :n] = 1.0

b = torch.zeros((n + 1,), device=A.device, dtype=A.dtype)
b[:n] = -g
b[n] = Q
return torch.linalg.solve(K, b)[:n]
\end{lstlisting}
The final element of the full solution is $\lambda$; slicing with
\texttt{[:n]} returns only the atomic charges.

This is a dense direct linear solve, not an iterative self-consistency loop.
Consequently, there is no QEq convergence tolerance or maximum iteration count.
The mathematical stationary point is obtained up to floating-point roundoff
and amplification by the condition number of $K$.  The augmented matrix is
indefinite because of the constraint row and column, but it is solvable when the
constrained QEq problem is well posed.

\section{Vacuum and polarized passes}

Each structure in a batch is solved separately.  The vacuum and polarized calls
are
\begin{lstlisting}[style=pythonstyle]
A = self._build_qeq_matrix(Rm, sigmam, Jiim)
q0m = self._solve_qeq(A, chim, Qm)
qphim = (
    self._solve_qeq(A, chim + phim, Qm)
    if polarize
    else q0m
)
\end{lstlisting}
Thus, $\vect q_0$ solves
\begin{equation}
  A\vect q_0+\vect\chi+\lambda_0\ones=0,
  \qquad \ones^{\mathsf T}\vect q_0=Q,
\end{equation}
whereas $\vect q_\phi$ solves
\begin{equation}
  A\vect q_\phi+\vect\chi+\vect\phi+\lambda_\phi\ones=0,
  \qquad \ones^{\mathsf T}\vect q_\phi=Q.
\end{equation}
The matrix $A$ is identical in the two solves; only the linear term changes.
Adding the same constant to all components of $\vect\phi$ does not change the
charge distribution under an exact constrained solve.  It changes only the
Lagrange multiplier.  This gauge invariance motivates molecule-wise centering
of fields in the training augmentation described in Appendix~\ref{app:field}.

\section{Finite-precision charge correction and warning}

Although Eq.~\eqref{eq:augmented-system} contains the total-charge constraint,
finite-precision arithmetic can leave a small residual.  For each structure the
implementation computes
\begin{equation}
  \Delta q=\frac{Q-\sum_iq_i}{N}
\end{equation}
and shifts every charge uniformly,
\begin{equation}
  q_i\leftarrow q_i+\Delta q.
\end{equation}
The code is
\begin{lstlisting}[style=pythonstyle]
q_sum = snn.scatter_add(q, idx_m, dim_size=maxm)
dq = (Q - q_sum) / natoms.unsqueeze(-1)
q = q + dq[idx_m]
\end{lstlisting}
Even after this operation, the subsequent floating-point summation can retain a
small residual.  The code warns if either
\begin{equation}
  r_0=\max_m\left|\sum_{i\in m}q_{0,i}-Q_m\right|
  \quad\text{or}\quad
  r_\phi=\max_m\left|\sum_{i\in m}q_{\phi,i}-Q_m\right|
\end{equation}
exceeds $10^{-5}e$.  This value is a post-solve diagnostic threshold, not a
solver convergence criterion.  A large or repeated residual can indicate an
ill-conditioned QEq matrix, extreme charge cancellation, or insufficient
floating-point precision.  The configured training precision is 32 bit.

\section{External-field embedding energy}

The implementation reports
\begin{equation}
  E_{\mathrm{emb}}
  =\frac{1}{2}(\vect q_0+\vect q_\phi)^{\mathsf T}\vect\phi,
  \label{eq:embedding-energy}
\end{equation}
using
\begin{lstlisting}[style=pythonstyle]
embedding_atom = 0.5 * (q0 + qphi) * phi
embedding_energy = scatter_add(embedding_atom, idx_m)
\end{lstlisting}
Equation~\eqref{eq:embedding-energy} is the linear-response charging expression:
as the external potential is turned on, the charge changes linearly from
$\vect q_0$ to $\vect q_\phi$, whose path average is
$(\vect q_0+\vect q_\phi)/2$.

\section{Core and valence diagnostic outputs}

Only after solving QEq does the implementation look up an element-dependent
core charge,
\begin{align}
  q_{\mathrm{val},i}^{\phi}&=q_{\phi,i}-q_{\mathrm{core},i},\\
  q_{\mathrm{val},i}^{0}&=q_{0,i}-q_{\mathrm{core},i}.
\end{align}
The corresponding code is
\begin{lstlisting}[style=pythonstyle]
q_core = self._get_qcore(Z).unsqueeze(-1)
q_val = qphi - q_core
q_val_vac = q0 - q_core
\end{lstlisting}
Therefore, \texttt{qcore\_table} does not enter the QEq matrix, linear term,
total-charge constraint, charge solution, or embedding energy.  It affects only
the reported core/valence decomposition.  In the present configuration the
table contains $1.0$ for H, $6.35$ for O, and $4.46$ for Fe.

\section{Computational characteristics}

For $N$ QEq atoms, explicitly storing the dense matrix costs
$\mathcal O(N^2)$ memory and its direct factorization costs approximately
$\mathcal O(N^3)$ work.  This approach is attractive for small and medium-sized
QEq regions because it has no iterative convergence error and is directly
differentiable through PyTorch.  For very large systems, a sparse or
matrix-free iterative method can be faster and use less memory.  A direct solve
is not automatically immune to numerical error: poor conditioning can amplify
roundoff in either direct or iterative methods.

\appendix
\section{Field-regularized training task}
\label{app:field}

\subsection{Purpose and two-pass training structure}

\texttt{FieldRegularizedAtomisticTask} augments ordinary supervised training
with a second, unlabeled external-field pass.  The purpose is to discourage
unphysical charge magnitudes and excessive field-induced charge transfer when
field-dependent reference labels are unavailable.

For every training batch, the task performs:
\begin{enumerate}
  \item a zero-field forward pass, used for the ordinary supervised loss; and
  \item an augmented-field forward pass, used only for charge regularization.
\end{enumerate}
When field regularization is disabled, the class delegates directly to the
standard \texttt{AtomisticTask}.  The override is confined to
\texttt{training\_step}; validation and testing follow the normal task path and
do not receive random field augmentation.

\subsection{Supervised zero-field pass}

The task makes a shallow copy of the batch and explicitly sets the external
potential to zero at every atom:
\begin{lstlisting}[style=pythonstyle]
zero_batch = dict(batch)
zero_batch[self.phi_key] = batch[properties.Z].new_zeros(
    batch[properties.Z].shape, dtype=batch[properties.R].dtype
)
pred_zero = self.predict_without_postprocessing(zero_batch)
q_zero = pred_zero[self.charges_key]
supervised_loss = self.loss_fn(pred_zero, targets)
\end{lstlisting}
For the supplied training configuration, the supervised composite loss is
\begin{equation}
  \mathcal L_{\mathrm{sup}}
  =0.05\,\mathrm{MSE}(V)
  +0.70\,\mathrm{MSE}(\vect F)
  +0.15\,\mathrm{MSE}(\vect q_0)
  +0.10\,\mathrm{MSE}(\vect s).
  \label{eq:supervised-loss}
\end{equation}
Here $V$ is \texttt{V\_BuRNN}, $\vect F$ is \texttt{F\_BuRNN}, and $\vect s$
denotes \texttt{mbis\_valence\_widths}.  The implementation stores
\texttt{q\_zero} before output constraints are applied, because such constraints
may replace entries in the prediction dictionary.

\subsection{Generation of the augmented field}

Two field modes are supported:
\begin{description}
  \item[\texttt{random}] Draw an independent value for each atom,
    $\widetilde\phi_i\sim\mathcal N(0,\sigma_\phi^2)$.
  \item[\texttt{batch}] Read the field from \texttt{phi\_static\_aug}, falling
    back to \texttt{phi\_static}; if neither is present, generate a random field.
\end{description}
In random mode the implementation starts from
\begin{lstlisting}[style=pythonstyle]
source = torch.randn_like(batch[properties.R][:, 0]) * self.phi_std
\end{lstlisting}
If molecule-wise centering is enabled, it applies
\begin{equation}
  \phi_i\leftarrow\widetilde\phi_i
  -\frac{1}{N_m}\sum_{j\in m}\widetilde\phi_j,
  \qquad i\in m.
  \label{eq:field-centering}
\end{equation}
This removes the physically redundant constant-potential component separately
for every molecule.  The result is then clipped component-wise:
\begin{equation}
  \phi_i\leftarrow\operatorname{clip}
  (\phi_i,-\phi_{\max},\phi_{\max}).
\end{equation}
Clipping after centering can reintroduce a small nonzero molecular mean when
individual components hit the bounds; this has no effect on exact constrained
QEq charge differences, although finite-precision effects remain possible.

The current configuration uses
\begin{equation}
  \sigma_\phi=300\ \mathrm{kJ\,mol^{-1}\,e^{-1}},
  \qquad
  \phi_{\max}=1000\ \mathrm{kJ\,mol^{-1}\,e^{-1}},
\end{equation}
with centering enabled.

\subsection{Auxiliary polarized pass}

The field is placed under the same key read by \texttt{EMLEQEqStatic}, and the
model is evaluated again:
\begin{lstlisting}[style=pythonstyle]
aug_batch = dict(batch)
aug_batch[properties.R] = batch[properties.R].detach().clone()
aug_batch[self.phi_key] = phi_aug
pred_aug = self.predict_without_postprocessing(aug_batch)
q_phi = pred_aug[self.charges_key]
\end{lstlisting}
The detached copy of the positions keeps force derivatives in this auxiliary
forward pass independent of the supervised force graph.  No auxiliary energy
or force label is used.  Gradients from the charge regularizers still propagate
through $\vect q_\phi$ into the QEq parameter head and the shared neural-network
representation.

\subsection{Element-dependent charge-bound penalty}

Let $b_i=b(Z_i)$ be the configured absolute charge bound for atom $i$.  The
penalty is a squared hinge loss,
\begin{equation}
  \mathcal L_{\mathrm{bound}}
  =\frac{1}{N_{\mathrm{batch\ atoms}}}
   \sum_i\left[\max\left(|q_{\phi,i}|-b_i,0\right)\right]^2.
  \label{eq:bound-loss}
\end{equation}
It is zero inside the allowed interval and grows quadratically outside it.  The
configured bounds are
\begin{center}
\begin{tabular}{@{}lcc@{}}
\toprule
Element & Atomic number & $b(Z)$ in $e$\\
\midrule
H  & 1  & 0.8\\
O  & 8  & 1.5\\
Fe & 26 & 3.2\\
Other elements & --- & 2.0\\
\bottomrule
\end{tabular}
\end{center}
The bound is on the total QEq charge $q_{\phi,i}$, not on the diagnostic valence
charge derived from \texttt{qcore\_table}.

\subsection{Field-induced charge-change penalty}

The second auxiliary term is
\begin{equation}
  \mathcal L_{\Delta q}
  =\frac{1}{N_{\mathrm{batch\ atoms}}}
   \sum_i(q_{\phi,i}-\operatorname{stopgrad}(q_{0,i}))^2.
  \label{eq:dq-loss}
\end{equation}
The code uses
\begin{lstlisting}[style=pythonstyle]
dq_loss = (q_phi - q_zero.detach()).square().mean()
\end{lstlisting}
Detaching $\vect q_0$ means this auxiliary term treats the zero-field charges as
a fixed reference during backpropagation.  It pushes only the polarized branch
toward that reference and does not use the regularizer to move the reference
charges themselves.  The term limits field sensitivity, but it is not an
explicit Jacobian norm such as $\lVert\partial\vect q/\partial\vect\phi\rVert^2$.
The configuration key \texttt{w\_sensitivity} is reserved for such a future
penalty and must currently be zero.

\subsection{Complete training objective}

The final per-batch objective is
\begin{equation}
  \boxed{
  \mathcal L
  =\mathcal L_{\mathrm{sup}}
  +w_{\mathrm{bound}}\mathcal L_{\mathrm{bound}}
  +w_{\Delta q}\mathcal L_{\Delta q}}
  \label{eq:total-loss}
\end{equation}
with the present weights
\begin{equation}
  w_{\mathrm{bound}}=0.01,
  \qquad
  w_{\Delta q}=0.001.
\end{equation}
Both regularizers are unsupervised: they require neither polarized reference
charges nor field-dependent energies and forces.  Their numerical influence
depends on their raw scales as well as the displayed weights, so logged loss
components should be inspected rather than interpreting the weights alone.

\subsection{Logged diagnostics}

During training, the task logs the following batch-level diagnostics:
\begin{itemize}
  \item maximum polarized absolute charge,
        \texttt{train\_q\_phi\_abs\_max};
  \item unweighted bound penalty,
        \texttt{train\_q\_phi\_bound\_loss};
  \item unweighted charge-change penalty,
        \texttt{train\_q\_phi\_dq\_loss}; and
  \item minimum, maximum, and population standard deviation of the augmented
        potential.
\end{itemize}
These quantities make it possible to distinguish an overly strong field
distribution from an overly responsive QEq model and to assess the actual scale
of each auxiliary contribution to Eq.~\eqref{eq:total-loss}.

\subsection{Current configuration summary}

\begin{center}
\begin{tabular}{@{}ll@{}}
\toprule
Setting & Current value\\
\midrule
Task & \texttt{FieldRegularizedAtomisticTask}\\
Field key read by QEq & \texttt{phi\_static}\\
Augmentation mode & \texttt{random}\\
Random-field standard deviation & $300\ \mathrm{kJ\,mol^{-1}\,e^{-1}}$\\
Absolute clip & $1000\ \mathrm{kJ\,mol^{-1}\,e^{-1}}$\\
Molecule-wise centering & enabled\\
$w_{\mathrm{bound}}$ & $0.01$\\
$w_{\Delta q}$ & $0.001$\\
Explicit Jacobian penalty & not implemented; weight fixed to zero\\
Training precision & 32 bit\\
\bottomrule
\end{tabular}
\end{center}

\end{document}
